% Part: first-order-logic
% Chapter: syntax
% Section: intro-syntax

\documentclass[../../../include/open-logic-section]{subfiles}

\begin{document}

\olfileid{fol}{syn}{itx}

\olsection{Inleiding}

Om de theorie en metalogica van de eerste-ordelogica te ontwikkelen,
moeten we eerst de syntaxis en semantiek van haar uitdrukkingen
definiëren. De uitdrukkingen van de eerste-ordelogica zijn termen en
!!{formula}s. Termen worden gevormd uit !!{variable}s, !!{constant}s en
!!{function}s. !!^{formula}s worden op hun beurt gevormd uit
!!{predicate}s samen met termen (deze vormen de kleinste, ``atomaire''
!!{formula}s); uit atomaire !!{formula}s kunnen we vervolgens met
behulp van logische connectieven en kwantoren complexere formules
vormen. Er zijn veel verschillende manieren om de vormingsregels vast
te leggen; wij geven slechts één mogelijkheid. Andere systemen kiezen
andere symbolen, nemen andere verzamelingen connectieven als primitief
en gebruiken haakjes op een andere manier (of helemaal niet, zoals in
de zogenoemde Poolse notatie). Alle benaderingen hebben echter gemeen
dat de vormingsregels de verzameling termen en !!{formula}s
\emph{inductief} definiëren. Als dit naar behoren gebeurt, kan iedere
uitdrukking volgens de vormingsregels in wezen slechts op één manier
ontstaan. Doordat de inductieve definitie \emph{eenduidig leesbare}
uitdrukkingen oplevert, kunnen we op dezelfde wijze betekenissen aan
die uitdrukkingen toekennen: door een inductieve definitie.
\end{document}
